% TODO make hw stylesheet
\documentclass[10pt]{article}
\usepackage[letterpaper,top=1in]{geometry}
\usepackage{quiver}
\usepackage[dvipsnames,svgnames]{xcolor}
\usepackage[shortlabels]{enumitem}
\usepackage[framemethod=TikZ]{mdframed}
\usepackage{amsmath,amssymb,amsthm}
\usepackage[colorlinks]{hyperref}
\usepackage{multicol}
\usepackage{tabularx}
\usepackage{stmaryrd}

\hypersetup{
	colorlinks=true,
	linkcolor=blue,
	filecolor=magenta,
	urlcolor=magenta,
	pdftitle={MAT 534 HW}
}

\usepackage{mathtools}
\usepackage{thmtools}
\usepackage[textsize=footnotesize]{todonotes}
\usepackage{titling}

\reversemarginpar
\mdfdefinestyle{mdbluebox}{roundcorner=10pt,innerbottommargin=9pt,
	linecolor=blue,backgroundcolor=TealBlue!5,}
\declaretheoremstyle[headfont=\sffamily\bfseries\color{MidnightBlue},
mdframed={style=mdbluebox},]{thmbluebox}

\mdfdefinestyle{mdbloobox}{frametitlefont=\bfseries,innerbottommargin=8pt,
	nobreak=true,backgroundcolor=TealBlue!5,linecolor=blue,}
\declaretheoremstyle[headfont=\bfseries\color{MidnightBlue},
mdframed={style=mdbloobox},headpunct={\\[3pt]},postheadspace=0pt,]{thmbloobox}

\mdfdefinestyle{mdredbox}{frametitlefont=\bfseries,innerbottommargin=8pt,
	nobreak=true,backgroundcolor=Salmon!5,linecolor=RawSienna,}
\declaretheoremstyle[headfont=\bfseries\color{RawSienna},
mdframed={style=mdredbox},headpunct={\\[3pt]},postheadspace=0pt,]{thmredbox}

\mdfdefinestyle{mdgreenbox}{linecolor=ForestGreen,backgroundcolor=ForestGreen!5,
	linewidth=2pt,rightline=false,leftline=true,topline=false,bottomline=false,}
\declaretheoremstyle[headfont=\bfseries\sffamily\color{ForestGreen!70!black},
mdframed={style=mdgreenbox},headpunct={ --- },]{thmgreenbox}

\mdfdefinestyle{mdorangebox}{linecolor=RedOrange,backgroundcolor=RedOrange!5,
	linewidth=2pt,rightline=false,leftline=true,topline=false,bottomline=false,}
\declaretheoremstyle[headfont=\bfseries\sffamily\color{RedOrange!70!black},
mdframed={style=mdorangebox},headpunct={ --- },]{thmorangebox}

\mdfdefinestyle{mdblackbox}{linecolor=black,backgroundcolor=RedViolet!5!gray!5,
	linewidth=3pt,rightline=false,leftline=true,topline=false,bottomline=false,}
\declaretheoremstyle[mdframed={style=mdblackbox}]{thmblackbox}

\declaretheoremstyle[spaceabove=6pt,spacebelow=6pt]{basehead}

\declaretheorem[style=basehead,name=Problem]{problem}
\declaretheorem[style=thmredbox,name=Problem,numbered=no]{problem*}
\declaretheorem[style=thmbloobox,name=Setup]{setup} % for config geo
\declaretheorem[style=thmredbox,name=Example,sibling=problem]{example}
\declaretheorem[style=thmredbox,name=$\bigstar$ Problem,sibling=problem]{niceprob}
\declaretheorem[style=thmbluebox,name=Theorem,sibling=problem]{theorem}
\declaretheorem[style=thmbluebox,name=Corollary,sibling=theorem]{corollary}
\declaretheorem[style=thmbluebox,name=Proposition,sibling=theorem]{prop}
\declaretheorem[style=thmbluebox,name=Lemma,numberwithin=problem]{lemma}
\declaretheorem[style=thmbluebox,name=Theorem,numbered=no]{theorem*}
\declaretheorem[style=thmbluebox,name=Lemma,numbered=no]{lemma*}
\declaretheorem[style=thmgreenbox,name=Claim,numberwithin=problem]{claim}
\declaretheorem[style=thmgreenbox,name=Claim,numbered=no]{claim*}
\declaretheorem[style=thmblackbox,name=Remark,numberwithin=problem]{remark}
\declaretheorem[style=thmblackbox,name=Remark,numbered=no]{remark*}
\declaretheorem[style=thmgreenbox,name=Definition,sibling=theorem]{definition}
\declaretheorem[style=thmgreenbox,name=Definition,numbered=no]{definition*}

\providecommand{\ol}{\overline}
\providecommand{\quiv}{\equiv}
\providecommand{\ceil}[1]{\left\lceil #1 \right\rceil}
\providecommand{\floor}[1]{\left\lfloor #1 \right\rfloor}
\newcommand{\dd}{\mathrm{d}}
\providecommand{\ul}{\underline}
\providecommand{\eps}{\varepsilon}
\providecommand{\half}{\frac{1}{2}}
\providecommand{\CC}{\mathbb C}
\providecommand{\FF}{\mathbb F}
\providecommand{\NN}{\mathbb N}
\providecommand{\QQ}{\mathbb Q}
\providecommand{\RR}{\mathbb R}
\providecommand{\ZZ}{\mathbb Z}
\providecommand{\dg}{^\circ}
\providecommand{\ii}{\item}
\providecommand{\setdiff}{\smallsetminus}
\providecommand{\alert}[1]{{\sffamily\textbf{\textcolor{blue}{#1}}}}
\providecommand{\inv}{^{-1}}
\providecommand{\mb}[1]{\mathbf{#1}}

\newenvironment{soln}{\begin{proof}[Solution]}{\end{proof}}

\providecommand{\pow}{\operatorname{Pow}}
\providecommand{\aut}{\operatorname{Aut}}
\providecommand{\vocab}[1]{{\textbf{\textcolor{ForestGreen}{#1}}}}
\providecommand{\lcm}{\operatorname{lcm}}
\providecommand{\abs}[1]{\left|#1\right|}
\providecommand{\Ob}{\operatorname{Ob}}
\providecommand{\Hom}{\operatorname{Hom}}
\providecommand{\Aut}{\operatorname{Aut}}
\providecommand{\id}{\operatorname{id}}
\newcommand{\doubleparen}[1]{(\!(#1)\!)}

\usepackage{mathrsfs}

% place title at top right
\setlength{\droptitle}{-8ex}
\pretitle{\begin{flushright}\Large\bfseries}
\posttitle{\par\end{flushright}}
\preauthor{\begin{flushright}}
\postauthor{\end{flushright}}
\predate{\begin{flushright}}
\postdate{\end{flushright}}

\title{MAT 534 HW07}
\author{\large Aditya Pahuja}
\date{\large Due 10/29 (extended to 10/31)}
\begin{document}
\maketitle

\begin{problem}
    Let $F$ be a field. Prove that the polynomial ring $F[x]$ is a PID.
\end{problem}

\begin{soln}
    Let $I\subseteq F[x]$ be an ideal and $f(x)\in I$ be a polynomial
    whose degree is as small as possible. Then for any $g(x)\in I$,
    we can divide $g$ by $f$ to get $g(x) = f(x)q(x) + r(x)$
    for some $q(x),r(x)\in F[x]$ where $\deg r < \deg f$;
    since $r(x) = g(x) - f(x)q(x)$, the only way this makes sense at all
    is if $r$ is the zero polynomial, i.e. $g(x)$ is divisible by $f(x)$,
    which implies that $I = (f(x))$.
\end{soln}

\begin{problem}
    Let $R$ be a PID.
    \begin{enumerate}[(a)]
        \ii Show that any two elements $a,b\in R$ have a least common multiple
	$\lcm(a,b)$, which is unique up to multiplication by units.
	\ii Show that $\gcd(a,b)\lcm(a,b)$ equals $ab$, up to a unit.
    \end{enumerate}
\end{problem}

\begin{soln}
    (a) Let $m$ be a generator of the intersection $(a)\cap(b)$.
    Then, if $M$ is a common multiple of $a$ and $b$,
    by definition $M\in (a)$ and $M\in (b)$, so $M\in(a)\cap (b) = (m)$,
    i.e. $m$ divides any common multiple $M$ of $a$ and $b$,
    making it a least common multiple of $a$ and $b$.

    (b) If $g$ is a greatest common divisor of $a$ and $b$, then
    write $a = ga'$, $b = gb'$. Then,
    $g = as + bt$ for some $s,t\in R$ since $(g) = (a) + (b)$,
    so $g = g(a's + b't)$, implying that $a's + b't = 1$
    because we can cancel the $g$s due to $R$ being an integral domain.
    Thus $(a') + (b') = (1)$, so $(a')\cap(b') = (a')(b') = (a'b')$,
    meaning $\lcm(a',b') = a'b'$.

    Now, $(g)((a')\cap(b')) \subseteq (a)\cap(b)$ because
    if $a'u = b'v\in (a')\cap(b')$ for some $u,v\in R$, then
    $au = ga'u = gb'v = bv$ is in $(a)\cap (b)$.
    The reverse inclusion holds because
    if $au = bv\in (a)\cap(b)$ then $ga'u = gb'v$ implies $a'u = b'v$
    due to being in an integral domain, so that
    $a'u = b'v\in (a')\cap(b')$ and $ga'u = gb'v\in (g)((a')\cap(b'))$.

    Therefore, $\lcm(a,b)\gcd(a,b) = g\lcm(ga',gb') = g^2\lcm(a',b')
    = ga'\cdot gb' = ab$.

    (All $\gcd$ and $\lcm$ computations in this problem are up to a unit.)
\end{soln}

\begin{problem}
    Let $R$ be an integral domain in which every prime ideal is principal.
    \begin{enumerate}[(a)]
        \ii Show that if $R$ is not a PID,
	then there is an ideal $I$ that is not principal,
	and is maximal with respect to this property.
	\ii Since $I$ cannot be prime, there are elements $a,b\in R$
	with $a,b\notin I$ and $ab\in I$. Show that $(a) + I = (c)$
	for some $c\in R$.
	\ii Show that the ideal $J = \{r\in R : rc\in I\}$ is principal.
	\ii Conclude that $I$ itself must be principal, which is a contradiction.
    \end{enumerate}
\end{problem}

\begin{soln}
    (a) Let $\{I_\alpha\}_{\alpha\in A}$ be a chain of non-principal ideals of $R$;
    the union $I = \bigcup_{\alpha\in A}I_\alpha$ is also an ideal,
    and it is non-principal since if $I = (a)$ for some $a\in R$
    then $a\in I_\alpha$ for some $\alpha\in A$, so
    $(a)\subseteq I_\alpha\subseteq I = (a)$, hence $I_\alpha$ must be principal,
    which it isn't. Therefore, by Zorn's lemma,
    the collection of non-principal ideals of $R$
    has a maximal element with respect to inclusion.

    (b) Since $a\notin I$, $(a) + I$ strictly contains $I$,
    so by the maximality of $I$ among the non-principal ideals,
    $(a) + I$ cannot be non-principal, i.e. is principal.

    (c) $J$ contains $I$ because if $x\in I$ then $xc\in I$,
    since ideals absorb multiplication; additionally, $J$ contains $b$
    because $bc\in b((a) + I) = (ab) + bI\subseteq I$,
    using the fact that $ab\in I$, so $J$ strictly contains $I$,
    hence is principal.

    (d) Let $J = (d)$ for some $d\in R$.
    If $u\in I\subseteq J$ then $u = cy$ for some $y = dz\in J$,
    so $u\in (cd)$, hence $I\subseteq (cd)$.
    On the other hand $d\in J$ means $cd\in I$
    so $(cd)\subseteq I$, so $I = (cd)$ is principal.
\end{soln}

\begin{problem}
    Exhibit all the ideals in $F[x]/(p(x))$, where $F$ is a field
    and $p(x)$ is a polynomial in $F[x]$.
\end{problem}

\begin{soln}
    Ideals in $F[x]/(p(x))$ are in correspondence with ideals of $F[x]$
    containing $p(x)$, which are exactly the ideals of the form
    $(q(x))$ where $q(x)$ is a polynomial dividing $p(x)$,
    with the correspondence given by sending $(q(x))$ to
    its image under the projection map $\phi\colon F[x]\to F[x]/(p(x))$;
    this image is just the ideal generated by $\phi(q(x))$.
    One collection of representatives of the elements of this ideal is
    the polynomials that divide $p(x)$ and are divisible by $q(x)$.
\end{soln}

\begin{problem}
    Let $p\in\ZZ$ be a prime with $p\quiv3\pmod 4$.
    Prove that the quotient ring $\ZZ[i]/(p)$ is a field with $p^2$ elements.
\end{problem}

\begin{soln}
    Suppose $p\mid ab$ for $a,b\in\ZZ[i]$.
    Then $p^2 = N(p)\mid N(ab) = N(a)N(b)$ where $N(x+yi) = x^2 + y^2$ for $x,y\in\ZZ$.
    If $p$ divides $N(x+yi) =x^2 + y^2$ for some $x,y\in\ZZ$, then
    either $x$ and $y$ are both multiples of $p$, or
    \begin{equation*}
	(xy\inv)^2\equiv -1\pmod p,
    \end{equation*}
    i.e. $-1$ is a quadratic residue modulo $p$.
    However, the latter is impossible because $p$ is $3\pmod 4$ ---
    to see this, if $g$ is a generator of the cyclic group $(\ZZ/p\ZZ)^\times$
    then $g^{(p-1)/2} = -1$ and $(p-1)/2$ is odd,
    while all nonzero quadratic residues are expressible as $g^k$
    for some even $k\in\{1,\dots,p-1\}$; this means $p\mid x + yi$
    is equivalent to $p$ dividing both $x$ and $y$.
    Since $p^2\mid N(a)N(b)$, $p$ divides at least one of $N(a)$ and $N(b)$,
    so $p$ divides whichever of $a$ and $b$ has norm divisible by $p$;
    thus $p$ is prime in $\ZZ[i]$.

    Now, $\ZZ[i]/(p) = \{x + yi : x,y\in\{0,\dots,p-1\}\}$
    clearly has $p^2$ elements (noting that these elements
    all represent different cosets of $(p)$ because
    the difference between any two distinct ones has norm indivisible by $p$)
    and is an integral domain since $(p)$ is a prime ideal,
    as it is generated by a prime element. Finite integral domains
    are fields, so $\ZZ[i]/(p)$ is a field of order $p^2$.
\end{soln}

\begin{problem}
    Show that $(x,y)$ is not a principal ideal in $\QQ[x,y]$.
\end{problem}

\begin{soln}
    Since $\QQ$ is a field, $\QQ[x,y]$ is a UFD.
    $x$ and $y$ are irreducible in this ring because
    if $x = pq$ for $p,q\in\QQ[x,y]$,
    then one of $p$ and $q$ must be a constant polynomial
    by degree considerations, and similarly for $y$.
    Therefore any common divisor of both $x$ and $y$ is a unit,
    so if $(x,y)$ is principal then it must equal $(1)$;
    but this is obviously false, since $(x,y)$ doesn't contain
    any constant polynomials.
\end{soln}

\begin{problem}
    Suppose that $f(x)$ and $g(x)$ are two polynomials
    with rational coefficients whose product $f(x)g(x)$ has integer coefficients.
    Show that the product of any coefficient of $f(x)$
    with any coefficient of $g(x)$ is an integer.
\end{problem}

\begin{soln}
    Let $q_f$ and $q_g$ be the unique rational numbers
    such that $q_ff(x)$ and $q_gg(x)$ are both primitive.
    Then $q_ff(x)\cdot q_gg(x)$ is primitive, so
    $\frac 1{q_fq_g}$ is equal to the content of $f(x)g(x)$, which is an integer.
    If $q_ff(x) = a_nx^n + \cdots a_1x + a_0$ and
    $q_gg(x) = b_mx^m + \cdots + b_1x + b_0$ for $a_i,b_j\in\ZZ$, then
    the coefficients of $f$ are $a_i/q_f$ and the coefficients of $g$ are
    $b_j/q_g$, so the product of some two is
    \begin{equation*}
	\frac{a_ib_j}{q_fq_g} = a_ib_j\cdot\operatorname{cont}(f(x)g(x))\in\ZZ.\qedhere
    \end{equation*}
\end{soln}

\begin{problem}
    Let $F$ be a field. A subring $R\subseteq F$ is called
    a \vocab{valuation ring} if, for every nonzero $x\in F$,
    at least one of $x$ and $x\inv$ belongs to $R$.
    \begin{enumerate}[(a)]
        \ii Show that $R$ has a unique maximal ideal.
	\ii Show that the ideals of $R$ are totally ordered under inclusion.
    \end{enumerate}
\end{problem}

\begin{soln}
    It suffices to show (b), since a totally ordered set has
    at most one maximal element.
    Let $I$ and $J$ be two ideals of $R$ that are not comparable under inclusion,
    so $I\not\subseteq J$ and $J\not\subseteq I$,
    and let $a\in I\setminus J$, $b\in J\setminus I$ (note $a,b\ne 0$).
    Then one of $ab\inv,a\inv b\in F$ is in $R$,
    so  either $ab\inv\cdot b = a\in J$ or
    $ba\inv\cdot a = b\in I$, contradiction;
    hence the collection of ideals of $R$ is totally ordered.
\end{soln}










\end{document}
